%==============================================================
%  lecons/ensembles/construction_Z.tex — Construction de Z
%==============================================================

\lecontitle{Construction de \texorpdfstring{$\mathbb{Z}$}{Z}}{Théorie des ensembles — Extension par différences}

%=============================================================
%  § 1 — MOTIVATION
%=============================================================
\secnum{navyblue}{1}{Motivation — inverser l'addition}

\rubriq{Pourquoi étendre $\mathbb{N}$ ?}
Dans $\mathbb{N}$, l'équation $n + x = m$ n'a pas toujours de solution
(elle n'en a que si $n \leq m$). On veut construire le plus petit anneau
contenant $\mathbb{N}$ dans lequel l'addition est inversible — i.e.\ un
\emph{groupe abélien}.

\begin{tcolorbox}[thmbox]
  \textbf{Idée.}\enspace%
  \index{entiers relatifs!construction}
  Un entier relatif est la «\,différence $m - n$\,» de deux entiers naturels.
  On modélise cela par une paire $(m, n) \in \mathbb{N}^2$ représentant $m - n$,
  avec la relation d'équivalence
  \[
    (m,n) \sim (m', n') \;\iff\; m + n' = m' + n.
  \]
\end{tcolorbox}

\decsep{}

%=============================================================
%  § 2 — CONSTRUCTION FORMELLE
%=============================================================
\secnum{navyblue}{2}{Construction formelle}

\begin{tcolorbox}[thmbox]
  \textbf{Définition.}\enspace%
  On pose $\mathbb{Z} = (\mathbb{N} \times \mathbb{N})/{\sim}$.
  On note $[m,n]$ la classe de $(m,n)$.
  Les lois sont définies par :
  \begin{align*}
    [m,n] + [p,q] &= [m+p,\; n+q], \\
    [m,n] \times [p,q] &= [mp+nq,\; mq+np].
  \end{align*}
  L'élément neutre de $+$ est $[0,0]$, l'opposé de $[m,n]$ est $[n,m]$.
  L'injection $\iota : \mathbb{N} \to \mathbb{Z},\; n \mapsto [n,0]$,
  est un morphisme de semi-anneaux (injection compatible avec $+$ et $\times$ ;
  $\mathbb{N}$ n'étant pas un anneau, on ne parle pas ici d'homomorphisme d'anneaux).
\end{tcolorbox}

\begin{mdframed}[style=proofbar]
  \textit{Bien-fondé de $+$.} Si $(m,n)\sim(m',n')$ et $(p,q)\sim(p',q')$,
  alors $m+n'=m'+n$ et $p+q'=p'+q$. En additionnant :
  $(m+p)+(n'+q') = (m'+p')+(n+q)$, donc $[m+p,n+q] = [m'+p',n'+q']$. $\checkmark$

  \textit{Bien-fondé de $\times$.} Plus délicat. On fixe d'abord $[p,q]$ et l'on
  remplace $(m,n)$ par $(m',n')$ équivalent : en multipliant l'identité
  $m+n'=m'+n$ par $p$ puis par $q$, on obtient
  $(mp+nq)+(m'q+n'p) = (m'p+n'q)+(mq+np)$, c'est-à-dire
  $[mp+nq,\,mq+np]=[m'p+n'q,\,m'q+n'p]$.
  On conclut de même en faisant varier le second facteur.

  \textit{Groupe abélien.} Associativité et commutativité héritées de $\mathbb{N}$.
  Neutre $[0,0]$, opposé $[n,m]$ : $[m,n]+[n,m]=[m+n,n+m]=[0,0]$ (car $m+n=n+m$).

  \textit{Anneau.} La distributivité se vérifie par calcul direct. $\blacksquare$
\end{mdframed}

\decsep{}

%=============================================================
%  § 3 — STRUCTURE ET PROPRIÉTÉS
%=============================================================
\secnum{navyblue}{3}{Structure de $\mathbb{Z}$}

\begin{tcolorbox}[thmbox]
  \textbf{Propriétés de $\mathbb{Z}$.}\enspace%
  \index{entiers relatifs!propriétés}
  \begin{enumerate}
    \item $(\mathbb{Z}, +, \times)$ est un \emph{anneau commutatif intègre}.
    \item $\mathbb{Z}$ est un \emph{anneau principal} : tout idéal de $\mathbb{Z}$
          est de la forme $n\mathbb{Z}$.
    \item L'ordre de $\mathbb{N}$ s'étend : $[m,n] \leq [p,q]
          \Leftrightarrow m+q \leq p+n$ dans $\mathbb{N}$.
          $(\mathbb{Z}, \leq)$ est un ordre total compatible avec $+$.
    \item $\mathrm{Frac}(\mathbb{Z}) = \mathbb{Q}$ (corps des fractions).
  \end{enumerate}
\end{tcolorbox}

\rubriq{Propriétés universelles}

Il faut distinguer deux énoncés de nature différente.

\begin{tcolorbox}[thmbox]
  \textbf{(a) Symétrisation : prolongement des morphismes additifs.}\enspace%
  Soit $(A,+)$ un groupe abélien et $\varphi : \mathbb{N} \to A$ un morphisme
  de monoïdes (i.e.\ $\varphi(m+n)=\varphi(m)+\varphi(n)$ et $\varphi(0)=0_A$).
  Il existe un unique morphisme de groupes $\tilde\varphi : \mathbb{Z} \to A$
  prolongeant $\varphi$, donné par $\tilde\varphi([m,n]) = \varphi(m) - \varphi(n)$.
  C'est la propriété universelle du \emph{groupe des différences} (symétrisation
  du monoïde $(\mathbb{N},+)$) ; elle vaut aussi pour les morphismes de
  semi-anneaux $\mathbb{N}\to A$ vers un anneau $A$, prolongés en morphismes
  d'anneaux.

  \smallskip
  \textbf{(b) Initialité de $\mathbb{Z}$ dans les anneaux unitaires.}\enspace%
  Pour tout anneau unitaire $A$, il existe un \emph{unique} homomorphisme
  d'anneaux unitaires $\chi : \mathbb{Z} \to A$. Autrement dit, $\mathbb{Z}$ est
  l'objet initial de la catégorie des anneaux unitaires.
\end{tcolorbox}

\begin{mdframed}[style=proofbar]
  \textit{Preuve de (b).} Un homomorphisme d'anneaux unitaires $\chi$ doit
  vérifier $\chi(1)=1_A$. Par additivité, $\chi(n)=n\cdot 1_A$ pour tout
  $n\in\mathbb{N}$, puis $\chi(-n)=-(n\cdot 1_A)$ : la valeur de $\chi$ sur tout
  $\mathbb{Z}$ est donc entièrement déterminée par $\chi(1)=1_A$, d'où
  l'unicité. Réciproquement, l'application $\chi(k)=k\cdot 1_A$ est bien définie
  (compatible avec $\sim$) et l'on vérifie que c'est un homomorphisme d'anneaux
  unitaires, d'où l'existence. $\blacksquare$
\end{mdframed}

\decsep{}

%=============================================================
%  § 4 — EXEMPLES, PIÈGES ET EXERCICES
%=============================================================
\secnum{exgreen}{4}{Exemples, pièges et exercices}

\rubriq{Exemples}

\begin{mdframed}[style=exbar]
  \textbf{1. Représentation standard.}\enspace%
  Tout $[m,n]$ est représenté soit par $[k,0]$ (entier $k = m-n \geq 0$)
  soit par $[0,k]$ avec $k \geq 1$ (entier $-k < 0$). On retrouve la
  présentation habituelle.

  \smallskip\noindent
  \textbf{2. Arithmétique.}\enspace%
  $\mathbb{Z}$ est euclidien (division euclidienne), donc principal, donc
  factoriel (décomposition en premiers unique à ordre et signe près).

  \smallskip\noindent
  \textbf{3. Caractéristique.}\enspace%
  $\mathbb{Z}$ est de caractéristique $0$ : $n \cdot 1 = [n,0] \neq [0,0]$
  pour tout $n \geq 1$.
\end{mdframed}

\vspace{6pt}
\rubriq{Pièges}

\begin{mdframed}[style=cexbar]
  \textbf{$\mathbb{Z}$ n'est pas un corps.}\enspace%
  $2 = [2,0]$ n'est pas inversible dans $\mathbb{Z}$ : $[2,0]\cdot[p,q]=[2p,2q]$,
  et $[2p,2q]=[1,0]$ imposerait $2p = 1+2q$, i.e.\ $2(p-q)=1$, impossible
  (membre de gauche pair, membre de droite impair).

  \smallskip\noindent
  \textbf{Divisibilité dans $\mathbb{Z}$ vs $\mathbb{N}$.}\enspace%
  Dans $\mathbb{Z}$, les unités sont $\pm 1$. Les irréductibles (ou premiers)
  se présentent par paires $\{p, -p\}$. Dans $\mathbb{N}$, seuls les $p > 0$
  sont considérés.
\end{mdframed}

\vspace{6pt}
\exolabel

\begin{mdframed}[style=exobar]
  \begin{enumerate}
    \item Montrer que la relation $\sim$ sur $\mathbb{N}^2$ est bien une relation
          d'équivalence.

    \item Montrer que $\mathbb{Z}$ est intègre : si $[m,n]\cdot[p,q]=[0,0]$
          et $[m,n]\neq[0,0]$, alors $[p,q]=[0,0]$.

    \item Montrer que tout idéal de $\mathbb{Z}$ est principal.
          (On utilisera la propriété du bon ordre de $\mathbb{N}$.)

    \item Montrer que $\mathbb{Z}/n\mathbb{Z}$ est un corps si et seulement si
          $n$ est premier.

    \item Soit $\sigma : \mathbb{Z} \to \mathbb{Z}$ un automorphisme d'anneau.
          Montrer que $\sigma = \mathrm{id}$.
  \end{enumerate}
\end{mdframed}

\leconfooter{Construction par paires due à Hamilton, Dedekind et Peano (fin \textsc{xix}\textsuperscript{e} s.) — formalisée dans ZFC au \textsc{xx}\textsuperscript{e} siècle}
